\documentclass[11pt,a4paper]{article}
\usepackage[margin=25mm]{geometry}
\usepackage[T1]{fontenc}
\usepackage[utf8]{inputenc}
\usepackage{lmodern,microtype,amsmath,amssymb,amsthm,mathtools}
\usepackage{booktabs,graphicx,float,hyperref,enumitem}
\hypersetup{colorlinks=true,linkcolor=blue,urlcolor=blue,citecolor=blue}
\numberwithin{equation}{section}
\newcommand{\dd}{\,\mathrm{d}}
\newcommand{\ii}{\mathrm{i}}
\newcommand{\e}{\mathrm{e}}
\newcommand{\E}[1]{\left\langle #1\right\rangle}
\newcommand{\problem}[2]{\subsection*{Problem #1: #2}\addcontentsline{toc}{subsection}{Problem #1: #2}}
\newcommand{\note}[1]{\par\medskip\noindent\textit{Check / interpretation.} #1\par}
\title{\textbf{The Wave Function}\\[2mm]\large Complete Worked Solutions to Chapter 1, Problems 1.1--1.18\\[1mm]\normalsize Independent Mathematical-Physics Study Notes}
\author{Kumar Azad\\\small \url{https://github.com/KrAzad0/QM-CH1-SOLUTION}}
\date{17 November 2026}
\begin{document}
\maketitle
\begin{abstract}
Original step-by-step calculations for all eighteen problems in Chapter~1 of D.~J.~Griffiths and D.~F.~Schroeter, \emph{Introduction to Quantum Mechanics}, third edition (Cambridge University Press, 2018). The emphasis is on normalization, expectation values, uncertainty, Ehrenfest's theorem, probability flux, classical comparison, and physical interpretation. Problem prompts are independent short paraphrases, not reproductions of the textbook. This is an \emph{arXiv-inspired layout}; it is not an arXiv submission or a peer-reviewed paper.
\end{abstract}
\tableofcontents
\section{Conventions and mathematical foundations}
For the one-dimensional Schrödinger equation with a real potential:
\begin{equation}
\ii\hbar\partial_t\Psi=\left[-\frac{\hbar^2}{2m}\partial_x^2+V(x,t)\right]\Psi,\quad
\rho=|\Psi|^2,\qquad \int_{\mathbb R}\rho\,\dd x=1.
\end{equation}
The key quantities for normalized states are
\begin{align}
\E{x^n}&=\int_{\mathbb R} x^n|\Psi|^2\dd x,&
\E{p^n}&=\int_{\mathbb R}\Psi^*(-\ii\hbar\partial_x)^n\Psi\dd x,\\
\sigma_A^2&=\E{A^2}-\E A^2,&
J&=\frac{\hbar}{m}\operatorname{Im}(\Psi^*\partial_x\Psi).
\end{align}
For a piecewise differentiable state, boundary terms require careful treatment. For exact textbook statements consult the reference below.
\section{Section exercises, Problems 1.1--1.9}
\problem{1.1}{Mean square and standard deviation for discrete ages}
\paragraph{Setup.} Fourteen ages occur with frequencies $(14,1),(15,1),(16,3),(22,2),(24,2),(25,5)$. Calculate the second moment, deviations, and variance.
\paragraph{Derivation.}
The total count is $N=1+1+3+2+2+5=14$ and $P(j)=N(j)/14$. The mean is
\begin{equation}
\E j=\frac{14+15+3(16)+2(22)+2(24)+5(25)}{14}
=\frac{294}{14}=21.
\end{equation}
The second moment and squared mean are
\begin{align}
\E{j^2}
&=\frac{14^2+15^2+3(16^2)+2(22^2)+2(24^2)+5(25^2)}{14}
=\frac{6434}{14}=\frac{3217}{7},\\
\E j^2&=441.
\end{align}
The deviations for the six distinct ages are $-7,-6,-5,1,3,4$. Directly,
\begin{align}
\sigma_j^2&=\frac{(-7)^2+(-6)^2+3(-5)^2+2(1)^2+2(3)^2+5(4)^2}{14}\\
&=\frac{260}{14}=\frac{130}{7}.
\end{align}
As a check, $\E{j^2}-\E j^2=3217/7-441=130/7$. Hence
\begin{equation}
\boxed{\sigma_j=\sqrt{\frac{130}{7}}\simeq4.30946.}
\end{equation}
\note{Both methods agree. An average of signed deviations would vanish and could not measure spread.}

\problem{1.2}{Falling-rock distribution: standard deviation and tails}
\paragraph{Setup.} For a uniformly timed snapshot during free fall through height $h$, the distance distribution is $\rho(x)=1/(2\sqrt{hx})$ for $0<x<h$. Find its standard deviation and probability outside one standard deviation from its mean.
\paragraph{Derivation.}
For $n>-1/2$ the moments are
\begin{equation}
\E{x^n}=\int_0^h\frac{x^n}{2\sqrt{hx}}\dd x
=\frac{1}{2\sqrt h}\left[\frac{x^{n+1/2}}{n+1/2}\right]_0^h
=\frac{h^n}{2n+1}.
\end{equation}
Thus $\mu=\E x=h/3$, $\E{x^2}=h^2/5$, and
\begin{equation}
\boxed{\sigma_x=\sqrt{\frac{h^2}5-\frac{h^2}9}=\frac{2h}{3\sqrt5}}.
\end{equation}
Both cutoffs $x_{\pm}=h(1/3\pm2/(3\sqrt5))$ lie inside $(0,h)$.
The cumulative distribution is $F(x)=\sqrt{x/h}$ on $[0,h]$. Hence
\begin{align}
P(|x-\mu|>\sigma_x)
&=F(x_-)+1-F(x_+)\\
&=\boxed{1+\sqrt{\frac13-\frac{2}{3\sqrt5}}-
\sqrt{\frac13+\frac{2}{3\sqrt5}}}
\simeq0.392938.
\end{align}
\note{The $x^{-1/2}$ singularity at the starting point is integrable. The usual Gaussian one-sigma heuristic does not apply to this distribution.}

\problem{1.3}{Normalization and moments of a Gaussian}
\paragraph{Setup.} Let $\rho(x)=A\exp[-\lambda(x-a)^2]$, with $\lambda>0$. Determine $A$, mean, second moment, variance, and shape.
\paragraph{Derivation.}
With $u=x-a$,
\begin{equation}
1=A\int_{-\infty}^{\infty}\e^{-\lambda u^2}\dd u
=A\sqrt{\frac{\pi}{\lambda}}
\quad\Rightarrow\quad
\boxed{A=\sqrt{\lambda/\pi}}.
\end{equation}
Oddness gives $\E u=0$ so $\E x=a$. Differentiating the Gaussian integral,
\begin{equation}
\int_{-\infty}^{\infty}u^2\e^{-\lambda u^2}\dd u
=-\frac{\dd}{\dd\lambda}\sqrt{\frac{\pi}{\lambda}}
=\frac{\sqrt\pi}{2\lambda^{3/2}}.
\end{equation}
Therefore $\E{u^2}=1/(2\lambda)$, and
\begin{equation}
\boxed{\E x=a,\qquad \E{x^2}=a^2+\frac1{2\lambda},
\qquad \sigma_x=\frac1{\sqrt{2\lambda}}.}
\end{equation}
The density is a symmetric bell with maximum $\sqrt{\lambda/\pi}$ at $x=a$, total area one, and inflection points at $a\pm1/\sqrt{2\lambda}$.

\problem{1.4}{Triangular wavefunction}
\paragraph{Setup.} Study the normalized piecewise state
\begin{equation}
\Psi(x,0)=
\begin{cases}
Ax/a,&0\le x\le a,\\
A(b-x)/(b-a),&a\le x\le b,\\
0,&\text{otherwise},
\end{cases}
\qquad 0<a<b.
\end{equation}
\paragraph{Derivation.}
The graph of $\Psi$ is triangular: it rises from zero at $0$ to $A$ at $a$ and falls to zero at $b$. Choosing the irrelevant overall phase so $A>0$,
\begin{align}
1&=A^2\left[\frac1{a^2}\int_0^a x^2\dd x+
\frac1{(b-a)^2}\int_a^b(b-x)^2\dd x\right]\\
&=A^2\left[\frac a3+\frac{b-a}3\right]
=A^2\frac b3,
\end{align}
and so $\boxed{A=\sqrt{3/b}}$. Since $\rho=|\Psi|^2$, the most likely position is the apex $x=a$ (the probability at an exact point is zero; this means the \emph{density} is maximal there).
The probability of finding $x<a$ is
\begin{equation}
P(x<a)=A^2\int_0^a \frac{x^2}{a^2}\dd x
=\boxed{\frac ab}.
\end{equation}
It tends to one as $b\downarrow a$, and equals $1/2$ when $b=2a$. For the mean,
\begin{align}
\E x&=A^2\left[\frac1{a^2}\int_0^a x^3\dd x+
\frac1{(b-a)^2}\int_a^b x(b-x)^2\dd x\right]\\
&=\frac3b\left[\frac{a^2}4+\frac{b(b-a)}3-\frac{(b-a)^2}4\right]\\
&=\boxed{\frac a2+\frac b4}.
\end{align}
\note{The mean of the amplitude's square is not the familiar $(a+b)/3$ mean of a triangular probability density. For $b=2a$, symmetry gives $\E x=a$, as required.}

\problem{1.5}{Two-sided exponential localization}
\paragraph{Setup.} The wavefunction is $\Psi=A\e^{-\lambda|x|}\e^{-\ii\omega t}$, with $\lambda>0$. Compute normalization, first two position moments, width, and probability beyond one standard deviation.
\paragraph{Derivation.}
The time factor has unit modulus, so
\begin{equation}
1=A^2\int_{\mathbb R}\e^{-2\lambda|x|}\dd x
=2A^2\int_0^\infty\e^{-2\lambda x}\dd x
=\frac{A^2}{\lambda}
\quad\Rightarrow\quad \boxed{A=\sqrt\lambda}.
\end{equation}
By parity $\E x=0$. Using $\int_0^\infty x^2\e^{-qx}\dd x=2/q^3$,
\begin{align}
\E{x^2}&=2\lambda\int_0^\infty x^2\e^{-2\lambda x}\dd x
=2\lambda\,\frac{2}{(2\lambda)^3}
=\frac1{2\lambda^2},\\
\boxed{\sigma_x}&=\boxed{\frac1{\sqrt2\,\lambda}}.
\end{align}
Finally
\begin{equation}
P(|x|>\sigma_x)
=2\lambda\int_{\sigma_x}^\infty \e^{-2\lambda x}\dd x
=\boxed{\e^{-\sqrt2}}\simeq0.243117.
\end{equation}
The graph of $|\Psi|^2=\lambda\e^{-2\lambda|x|}$ is a symmetric exponential cusp; mark the points $x=\pm1/(\sqrt2\lambda)$.

\problem{1.6}{Why a time derivative cannot be integrated by parts in space}
\paragraph{Setup.} Explain the mistaken inference $\dd\E x/\dd t=0$ obtained by trying to transfer $\partial_t$ onto $x$ under a spatial integral.
\paragraph{Derivation.}
The derivative
\begin{equation}
\frac{\dd\E x}{\dd t}=\int_{\mathbb R}x\,\partial_t|\Psi(x,t)|^2\dd x
\end{equation}
is with respect to $t$ whereas integration is with respect to $x$. The integration-by-parts formula
\begin{equation}
\int_a^b f(x)\,\partial_xg(x)\dd x=[fg]_a^b-\int_a^b g(x)\,\partial_x f(x)\dd x
\end{equation}
requires the derivative and integration variable to match. One may not change $\partial_t$ to $\partial_x$. Instead, after invoking the Schrödinger equation and continuity equation $\partial_t\rho=-\partial_xJ$,
\begin{equation}
\frac{\dd\E x}{\dd t}
=-\int x\,\partial_xJ\dd x
=-[xJ]_{-\infty}^{\infty}+\int J\dd x
=\frac{\E p}{m},
\end{equation}
assuming the boundary terms vanish. The last identity follows by integrating by parts in the current and using the definition of $\E p$.

\problem{1.7}{Ehrenfest's equation for momentum}
\paragraph{Setup.} Derive the time evolution of the momentum expectation for a particle in a real scalar potential $V(x,t)$.
\paragraph{Derivation.}
For a wavefunction satisfying the Schrödinger equation and the standard boundary conditions, any time-independent operator $\hat A$ obeys
\begin{equation}
\frac{\dd}{\dd t}\E A
=\frac{\ii}{\hbar}\E{[\hat H,\hat A]},
\qquad \hat H=\frac{\hat p^2}{2m}+V(x,t).
\end{equation}
This follows by differentiating $\int\Psi^*\hat A\Psi\,\dd x$ and substituting
$\partial_t\Psi=-(\ii/\hbar)\hat H\Psi$, together with its Hermitian conjugate. For $\hat A=\hat p=-\ii\hbar\partial_x$, the kinetic term commutes with $\hat p$, whereas
\begin{align}
[V,\hat p]\psi
&=V(-\ii\hbar\psi')-(-\ii\hbar\partial_x)(V\psi)\\
&=-\ii\hbar V\psi'+\ii\hbar(V'\psi+V\psi')
=\ii\hbar V'\psi.
\end{align}
Therefore
\begin{equation}
\boxed{\frac{\dd\E p}{\dd t}
=\frac{\ii}{\hbar}\E{[V,\hat p]}
=-\E{\frac{\partial V}{\partial x}}.}
\end{equation}
\note{This is the average-force form of Newton's second law. It does not generally mean $\dd\E p/\dd t=-V'(\E x)$, because averaging and evaluating a nonlinear function at the mean are different operations.}

\problem{1.8}{Adding a constant to potential energy}
\paragraph{Setup.} Compare the solutions of the Schrödinger equation for $V(x,t)$ and $V'(x,t)=V(x,t)+V_0$, where $V_0$ is real and independent of position and time.
\paragraph{Derivation.}
Suppose $\Psi$ solves $\ii\hbar\partial_t\Psi=\hat H\Psi$. Let
\begin{equation}
\Psi'(x,t)=\e^{-\ii V_0t/\hbar}\Psi(x,t).
\end{equation}
Differentiating,
\begin{align}
\ii\hbar\partial_t\Psi'
&=\e^{-\ii V_0t/\hbar}
\left[V_0\Psi+\ii\hbar\partial_t\Psi\right]\\
&=(\hat H+V_0)\Psi',
\end{align}
so $\Psi'$ is the desired state for the shifted potential. Since the extra factor is a spatially uniform complex phase,
\begin{equation}
|\Psi'|^2=|\Psi|^2,\qquad
\E{Q}'=\int\Psi'^*\hat Q\Psi'\dd x
=\int\Psi^*\hat Q\Psi\dd x=\E Q
\end{equation}
for physical observables $\hat Q$ unaffected by the shift. The Hamiltonian expectation itself increases by $V_0$ because the \emph{energy origin} has changed, even though position, momentum, and other unshifted observables do not change.
\note{A constant potential offset adds a global dynamical phase and shifts all energy levels equally. Energy differences and measurement probabilities remain unchanged.}

\problem{1.9}{A Gaussian oscillator state saturating uncertainty}
\paragraph{Setup.} Consider $\Psi(x,t)=A\exp[-a(mx^2/\hbar+\ii t)]$ with real $a>0$. Determine $A$, the required potential, the first two moments of position and momentum, and their uncertainties.
\paragraph{Derivation: normalization.}
Taking the absolute square eliminates the time phase,
\begin{equation}
|\Psi|^2=A^2\e^{-2amx^2/\hbar}.
\end{equation}
The Gaussian integral yields
\begin{equation}
1=A^2\sqrt{\frac{\pi\hbar}{2am}}
\quad\Rightarrow\quad
\boxed{A=\left(\frac{2am}{\pi\hbar}\right)^{1/4}}.
\end{equation}
\paragraph{Derivation: inverse Schrödinger equation.}
The derivatives are
\begin{align}
\partial_t\Psi&=-\ii a\Psi,&
\partial_x\Psi&=-\frac{2amx}{\hbar}\Psi,\\
\partial_x^2\Psi&=\left(-\frac{2am}{\hbar}+
\frac{4a^2m^2x^2}{\hbar^2}\right)\Psi.
\end{align}
Therefore
\begin{align}
\ii\hbar\partial_t\Psi&=\hbar a\Psi,\\
-\frac{\hbar^2}{2m}\partial_x^2\Psi
&=\left(\hbar a-2ma^2x^2\right)\Psi.
\end{align}
The Schrödinger equation holds precisely for
\begin{equation}
\boxed{V(x)=2ma^2x^2=\frac12m(2a)^2x^2,}
\end{equation}
which is a harmonic-oscillator potential of frequency $\omega=2a$, and the state has energy $\hbar a=\hbar\omega/2$.
\paragraph{Derivation: moments and widths.}
The spatial probability density is even, so $\E x=0$. Its Gaussian variance is
\begin{equation}
\boxed{\E{x^2}=\frac{\hbar}{4am}}.
\end{equation}
We likewise obtain
\begin{equation}
\E p=-\ii\hbar\int\Psi^*\partial_x\Psi\dd x
=2\ii am\E x=0,
\end{equation}
and either integrating $\Psi^*(-\hbar^2\partial_x^2)\Psi$ or using
$\hbar^2\int|\partial_x\Psi|^2$ gives
\begin{equation}
\boxed{\E{p^2}=4a^2m^2\E{x^2}=am\hbar}.
\end{equation}
Consequently
\begin{equation}
\boxed{\sigma_x=\sqrt{\frac{\hbar}{4am}},\qquad
\sigma_p=\sqrt{am\hbar},\qquad
\sigma_x\sigma_p=\frac{\hbar}{2}.}
\end{equation}
\note{Equality holds: this real Gaussian is the ground state of the oscillator. The parameter $a$ has dimensions of inverse time.}

\section{Further exercises, Problems 1.10--1.18}
\problem{1.10}{Empirical statistics of the first 25 digits of $\pi$}
\paragraph{Setup.} Treat the first twenty-five decimal-expansion digits, beginning with the initial $3$, as an equally likely sample space. Determine digit probabilities, mode, median, mean, and standard deviation.
\paragraph{Derivation.}
The sequence is
\begin{equation}
3,\ 1,\ 4,\ 1,\ 5,\ 9,\ 2,\ 6,\ 5,\ 3,\ 5,\ 8,\ 9,\
7,\ 9,\ 3,\ 2,\ 3,\ 8,\ 4,\ 6,\ 2,\ 6,\ 4,\ 3.
\end{equation}
Counting occurrences gives the following frequency table, with probability $P(j)=N(j)/25$:
\begin{center}
\begin{tabular}{c*{10}{r}}
\toprule
Digit $j$&0&1&2&3&4&5&6&7&8&9\\
\midrule
Count $N(j)$&0&2&3&5&3&3&3&1&2&3\\
\bottomrule
\end{tabular}
\end{center}
The total count is $0+2+3+5+3+3+3+1+2+3=25$. The digit $3$ occurs five times and is the unique mode. For twenty-five ordered observations, the median is the thirteenth; the cumulative counts through digits $0,1,2,3$ are $0,2,5,10$ and through digit $4$ are $13$, so the median is $4$. For the moments,
\begin{align}
\E j&=\frac{1(2)+2(3)+3(5)+4(3)+5(3)+6(3)+7(1)+8(2)+9(3)}{25}
=\boxed{\frac{118}{25}=4.72},\\
\E{j^2}&=\frac{1^2(2)+2^2(3)+3^2(5)+4^2(3)+5^2(3)+6^2(3)+7^2(1)+8^2(2)+9^2(3)}{25}
=\frac{710}{25}=28.4.
\end{align}
It follows that
\begin{equation}
\boxed{\sigma_j=\sqrt{28.4-(4.72)^2}=\sqrt{6.1216}\simeq2.47419.}
\end{equation}
\note{A frequency table can serve as a discrete probability distribution even when the sampled symbols are digits rather than physical measurements.}

\problem{1.11}{Classical position density of a harmonic oscillator}
\paragraph{Setup.} For a particle moving at fixed energy $E$ between classical turning points, obtain the time-sampling density from its speed, then specialize to $V(x)=kx^2/2$ and calculate its position moments.
\paragraph{Derivation: general energy.}
Energy conservation is
\begin{equation}
E=\frac12mv(x)^2+V(x)\quad\Rightarrow\quad
\boxed{v(x)=\sqrt{\frac2m[E-V(x)]}}.
\end{equation}
For the travel time $T=\int_a^b \dd x/v(x)$ between turning points, a random observation during the \emph{one-way} journey has normalized density
\begin{equation}
\rho(x)=\frac1{T\,v(x)},\qquad
\int_a^b\rho(x)\dd x=\frac1T\int_a^b\frac{\dd x}{v(x)}=1.
\end{equation}
Over a full periodic orbit, each interior interval is visited twice; the factor of two appears both in residence time and period and therefore cancels.
\paragraph{Derivation: harmonic potential.}
Let $\omega=\sqrt{k/m}$ and $A=\sqrt{2E/k}$, so the turning points are $\pm A$. Then
\begin{equation}
v(x)=\omega\sqrt{A^2-x^2},\qquad
T=\frac1\omega\int_{-A}^A\frac{\dd x}{\sqrt{A^2-x^2}}
=\frac{\pi}{\omega}.
\end{equation}
Thus
\begin{equation}
\boxed{\rho(x)=\frac1{\pi\sqrt{A^2-x^2}},\quad |x|<A;\qquad
\rho(x)=0\ \text{outside}.}
\end{equation}
It has integrable endpoint peaks at $\pm A$, where the particle moves most slowly. To compute moments set $x=A\sin\theta$ with $\theta\in[-\pi/2,\pi/2]$:
\begin{align}
\E x&=\frac A\pi\int_{-\pi/2}^{\pi/2}\sin\theta\,\dd\theta=0,\\
\E{x^2}&=\frac{A^2}{\pi}\int_{-\pi/2}^{\pi/2}\sin^2\theta\,\dd\theta
=\frac{A^2}{2},\\
\boxed{\sigma_x}&=\boxed{\frac A{\sqrt2}=\sqrt{\frac E{k}}}.
\end{align}
\note{The graph has a central minimum $1/(\pi A)$ and rises symmetrically toward the turning points. Classical position density differs strongly from the quantum ground-state Gaussian.}

\problem{1.12}{Classical momentum density and uncertainty product}
\paragraph{Setup.} Use the classical oscillator from Problem 1.11 to obtain momentum statistics and compare with the quantum uncertainty scale.
\paragraph{Derivation.}
For harmonic motion $x(t)=A\cos(\omega t)$ the momentum is
$p(t)=-mA\omega\sin(\omega t)$. Define
\begin{equation}
p_{\max}=mA\omega=\sqrt{2mE}.
\end{equation}
Uniform phase sampling, together with the two branches of the sine map in each period, gives the momentum density
\begin{equation}
\boxed{\rho_p(p)=\frac1{\pi\sqrt{p_{\max}^2-p^2}},
\quad |p|<p_{\max};\qquad 0\ \text{otherwise}.}
\end{equation}
Indeed, substituting $p=p_{\max}\sin\theta$ shows
$\int_{-p_{\max}}^{p_{\max}}\rho_p(p)\dd p=1$. By parity $\E p=0$ and
\begin{align}
\E{p^2}&=\frac{p_{\max}^2}{\pi}\int_{-\pi/2}^{\pi/2}\sin^2\theta\dd\theta
=\frac{p_{\max}^2}{2}=mE,\\
\boxed{\sigma_p}&=\boxed{\frac{p_{\max}}{\sqrt2}=\sqrt{mE}}.
\end{align}
Combining with $\sigma_x=A/\sqrt2$,
\begin{equation}
\boxed{\sigma_x\sigma_p=\frac{Ap_{\max}}2=\frac E\omega.}
\end{equation}
Classical mechanics permits $E\to0$, so the product can tend to zero. If one \emph{formally} inserts the oscillator's quantum lower bound $E\ge\hbar\omega/2$, this same expression yields
\begin{equation}
\sigma_x\sigma_p\ge\hbar/2.
\end{equation}
\note{The final inequality is an instructive comparison, not a derivation of quantum uncertainty from the classical distribution. The genuine quantum calculation must use a wavefunction and momentum operator.}

\problem{1.13}{Numerical verification by random phase sampling}
\paragraph{Setup.} Simulate ten thousand positions of a unit-amplitude oscillator sampled at random times and compare the histogram with the analytical density.
\paragraph{Analysis and algorithm.}
Set $A=\omega=1$. A uniform random time $t$ in $[0,2\pi]$ gives $X=\cos t$. For $-1<x<1$, the transformation has two roots per period, and
\begin{equation}
\rho_X(x)=\sum_{t_i:\cos t_i=x}\frac{1}{2\pi}\left|\frac{\dd t}{\dd x}\right|_{t_i}
=\frac{1}{\pi\sqrt{1-x^2}}.
\end{equation}
Its analytical mean and variance are $\E X=0$ and $\operatorname{Var}X=1/2$. A reproducible computation is:
\begin{enumerate}[leftmargin=*]
\item Draw $N=10\,000$ independent angles $\theta_j\sim U(0,2\pi)$.
\item Compute $x_j=\cos\theta_j$.
\item Plot a density-normalized histogram with e.g. $100$ bins spanning $[-1,1]$.
\item Superimpose $\rho(x)=1/[\pi\sqrt{1-x^2}]$ for $|x|<1$ (avoid evaluating exactly at the endpoints).
\item Compare the sample mean and sample standard deviation with $0$ and $1/\sqrt2\simeq0.7071$.
\end{enumerate}
The accompanying script \texttt{scripts/problem\_1\_13.py} produces a PDF/PNG plot and prints the sample statistics. The expected histogram is U-shaped, with large counts near the turning points; sampling fluctuations decrease roughly as $N^{-1/2}$ for fixed interior bins. The endpoint density diverges but each finite histogram bin has finite probability.
\IfFileExists{figures/problem_1_13.pdf}{
\begin{figure}[H]\centering
\includegraphics[width=.78\linewidth]{figures/problem_1_13.pdf}
\caption{Ten thousand random oscillator snapshots and the analytical arcsine density.}
\end{figure}}{\paragraph{Figure.} Run the companion Python script to generate \texttt{figures/problem\_1\_13.pdf} and recompile this manuscript to insert the comparison plot.}

\problem{1.14}{Continuity equation and probability current}
\paragraph{Setup.} For a normalized one-dimensional wavefunction, derive the time rate of probability in $(a,b)$ and compute the current for Problem 1.9.
\paragraph{Derivation: continuity equation.}
Let $\rho=\Psi^*\Psi$. For a real potential,
\begin{align}
\partial_t\Psi&=\frac{\ii\hbar}{2m}\partial_x^2\Psi
-\frac{\ii}{\hbar}V\Psi,\\
\partial_t\Psi^*&=-\frac{\ii\hbar}{2m}\partial_x^2\Psi^*
+\frac{\ii}{\hbar}V\Psi^*.
\end{align}
By the product rule, the potential terms cancel:
\begin{align}
\partial_t\rho
&=\Psi^*\partial_t\Psi+\Psi\,\partial_t\Psi^*\\
&=\frac{\ii\hbar}{2m}\left(
\Psi^*\partial_x^2\Psi-\Psi\partial_x^2\Psi^*\right)\\
&=\frac{\ii\hbar}{2m}\partial_x\left(
\Psi^*\partial_x\Psi-\Psi\partial_x\Psi^*\right)
=-\partial_xJ,
\end{align}
where
\begin{equation}
\boxed{J(x,t)=\frac{\hbar}{2m\ii}
\left(\Psi^*\partial_x\Psi-\Psi\partial_x\Psi^*\right)
=\frac{\hbar}{m}\operatorname{Im}(\Psi^*\partial_x\Psi).}
\end{equation}
Define $P_{ab}(t)=\int_a^b\rho(x,t)\dd x$. Then
\begin{equation}
\boxed{\frac{\dd P_{ab}}{\dd t}
=-\int_a^b\partial_xJ\dd x=J(a,t)-J(b,t).}
\end{equation}
The sign has an intuitive meaning: inflow through the left boundary increases probability; outflow through the right boundary decreases it.
\paragraph{Dimensions.}
In one spatial dimension $[\Psi]=L^{-1/2}$ and
$[\Psi^*\partial_x\Psi]=L^{-2}$. Since $[\hbar/m]=L^2/T$,
\begin{equation}
\boxed{[J]=T^{-1},\quad \text{probability per unit time}.}
\end{equation}
\paragraph{Application to Problem 1.9.}
There $\Psi(x,t)=f(x)\e^{-\ii at}$ with a real Gaussian $f(x)$, so
$\Psi^*\partial_x\Psi=f(x)f'(x)$ is real. Hence
\begin{equation}
\boxed{J(x,t)=0\quad\text{for every }x,t.}
\end{equation}
\note{A stationary probability distribution can still carry a current in other systems, but the Gaussian ground state considered here does not.}

\problem{1.15}{Conservation of inner products}
\paragraph{Setup.} Show that the overlap of any two normalizable solutions of the same Schrödinger equation with real $V(x,t)$ is independent of time.
\paragraph{Derivation.}
Define the overlap
\begin{equation}
S(t)=\int_{-\infty}^{\infty}\Psi_1^*(x,t)\Psi_2(x,t)\dd x
=\langle\Psi_1(t)|\Psi_2(t)\rangle.
\end{equation}
Using $\partial_t|\Psi_2\rangle=-\ii\hat H|\Psi_2\rangle/\hbar$ and
$\partial_t\langle\Psi_1|=\ii\langle\Psi_1|\hat H/\hbar$ for a self-adjoint Hamiltonian,
\begin{align}
\frac{\dd S}{\dd t}
&=\frac{\ii}{\hbar}\langle\Psi_1|\hat H|\Psi_2\rangle
-\frac{\ii}{\hbar}\langle\Psi_1|\hat H|\Psi_2\rangle\\
&=\boxed{0}.
\end{align}
To see the boundary assumption explicitly, the difference in the two Hamiltonian matrix elements is
\begin{align}
\langle H\Psi_1|\Psi_2\rangle-\langle\Psi_1|H\Psi_2\rangle
&=-\frac{\hbar^2}{2m}\int
\left[(\partial_x^2\Psi_1)^*\Psi_2-
\Psi_1^*\partial_x^2\Psi_2\right]\dd x\\
&=-\frac{\hbar^2}{2m}
\left[(\partial_x\Psi_1)^*\Psi_2-\Psi_1^*\partial_x\Psi_2
\right]_{-\infty}^{\infty}=0,
\end{align}
provided the boundary term vanishes.
Thus
\begin{equation}
\boxed{\langle\Psi_1(t)|\Psi_2(t)\rangle
=\langle\Psi_1(0)|\Psi_2(0)\rangle.}
\end{equation}
\note{Taking $\Psi_1=\Psi_2$ recovers norm conservation. Taking initially orthogonal states preserves their orthogonality under common unitary evolution.}

\problem{1.16}{Parabolic compact-support wavefunction}
\paragraph{Setup.} For $\Psi(x,0)=A(a^2-x^2)$ on $[-a,a]$ and zero outside, compute normalization, expectations $\E x,\E p,\E{x^2},\E{p^2}$, and uncertainty product.
\paragraph{Derivation: normalization.}
Choosing $A>0$,
\begin{align}
1&=A^2\int_{-a}^a(a^2-x^2)^2\dd x\\
&=A^2\int_{-a}^a(a^4-2a^2x^2+x^4)\dd x\\
&=A^2\left[2a^5-\frac{4a^5}3+\frac{2a^5}5\right]
=A^2\frac{16a^5}{15}.
\end{align}
Thus
\begin{equation}
\boxed{A=\frac{\sqrt{15}}{4a^{5/2}}.}
\end{equation}
\paragraph{Derivation: position moments.}
The density is even, giving $\E x=0$. For the second moment,
\begin{align}
\E{x^2}
&=A^2\int_{-a}^a x^2(a^2-x^2)^2\dd x\\
&=A^2\cdot2\left[\frac{a^7}3-\frac{2a^7}5+\frac{a^7}7\right]\\
&=\frac{15}{16a^5}\frac{16a^7}{105}
=\boxed{\frac{a^2}{7}}.
\end{align}
Therefore $\boxed{\sigma_x=a/\sqrt7}$.
\paragraph{Derivation: momentum moments.}
Inside the support, $\partial_x\Psi=-2Ax$; the wavefunction is real and vanishes at $x=\pm a$. Thus
\begin{equation}
\E p=-\ii\hbar\int_{-a}^a\Psi\,\partial_x\Psi\,\dd x
=-\frac{\ii\hbar}{2}\left[\Psi(x,0)^2\right]_{-a}^a
=\boxed{0}.
\end{equation}
Integration by parts, with zero endpoint contribution, yields
\begin{align}
\E{p^2}&=-\hbar^2\int_{-a}^a\Psi\,\partial_x^2\Psi\,\dd x
=\hbar^2\int_{-a}^a|\partial_x\Psi|^2\dd x\\
&=4A^2\hbar^2\int_{-a}^a x^2\dd x
=4\frac{15}{16a^5}\hbar^2\frac{2a^3}{3}
=\boxed{\frac{5\hbar^2}{2a^2}}.
\end{align}
As a result
\begin{equation}
\boxed{\sigma_p=\sqrt{\frac52}\frac{\hbar}{a}},\qquad
\boxed{\sigma_x\sigma_p=\hbar\sqrt{\frac5{14}}
\simeq0.597614\hbar>\frac{\hbar}{2}.}
\end{equation}
\note{Although the derivative jumps at the support endpoints, the wavefunction itself is continuous and vanishes there, and the weak derivative is square-integrable. The integration-by-parts quadratic form above is legitimate.}

\problem{1.17}{A complex potential as a decay model}
\paragraph{Setup.} Let a phenomenological potential have a constant absorptive imaginary part,
$V(x)=V_0(x)-\ii\Gamma$ with $\Gamma>0$. Determine the total probability and lifetime.
\paragraph{Derivation.}
Taking the Schrödinger equation and its complex conjugate as in Problem 1.14, the real part $V_0$ cancels, but the imaginary contribution produces a sink:
\begin{equation}
\partial_t|\Psi|^2+\partial_xJ
=-\frac{2\Gamma}{\hbar}|\Psi|^2.
\end{equation}
Integrating over all space and assuming no boundary flux,
\begin{align}
P(t)&=\int_{\mathbb R}|\Psi(x,t)|^2\dd x,\\
\frac{\dd P}{\dd t}&=-[J]_{-\infty}^{\infty}
-\frac{2\Gamma}{\hbar}P
=\boxed{-\frac{2\Gamma}{\hbar}P}.
\end{align}
Separate variables and integrate:
\begin{equation}
\frac{\dd P}{P}=-\frac{2\Gamma}{\hbar}\dd t
\quad\Rightarrow\quad P(t)=P(0)\e^{-2\Gamma t/\hbar}.
\end{equation}
For $P(0)=1$, comparison with $P(t)=\e^{-t/\tau}$ yields
\begin{equation}
\boxed{\tau=\frac{\hbar}{2\Gamma}}.
\end{equation}
\note{A non-Hermitian effective Hamiltonian does not conserve the norm of the selected subsystem. The decreasing norm represents survival probability and does not contradict probability conservation in a larger closed system.}

\problem{1.18}{Thermal de Broglie wavelength and quantum regimes}
\paragraph{Setup.} Estimate the threshold for thermal de Broglie wavelength $\lambda$ to exceed a characteristic separation $d$ in a solid or ideal gas. Compare electrons, silicon nuclei, helium gas, and dilute atomic hydrogen in space.
\paragraph{Derivation: general condition.}
The stated thermal kinetic-energy estimate $p^2/(2m)=3k_BT/2$ gives a typical momentum $p=\sqrt{3mk_BT}$ and
\begin{equation}
\boxed{\lambda=\frac{h}{\sqrt{3mk_BT}}}.
\end{equation}
The criterion $\lambda\gtrsim d$ becomes
\begin{equation}
\boxed{T\lesssim T_{\rm q}(m,d):=\frac{h^2}{3mk_Bd^2}.}
\end{equation}
This estimate indicates when \emph{translational} wave effects are important; it is a rule of thumb, not a sharp phase transition.
\paragraph{Solids.}
For $d=0.3\,\mathrm{nm}=3.0\times10^{-10}\,\mathrm{m}$, using $h=6.62607015\times10^{-34}\,\mathrm{J\,s}$ and $k_B=1.380649\times10^{-23}\,\mathrm{J/K}$,
\begin{align}
T_{\rm q,e}&=
\frac{h^2}{3(9.10938\times10^{-31}\,\mathrm{kg})k_Bd^2}
\simeq\boxed{1.29\times10^5\,\mathrm K},\\
T_{\rm q,Si}&=
\frac{h^2}{3(28.085\,u)k_Bd^2}
\simeq\boxed{2.53\,\mathrm K},\qquad
u=1.66053906660\times10^{-27}\,\mathrm{kg}.
\end{align}
Thus conducting electrons are typically quantum-mechanical even at room temperature by this criterion, whereas lattice nuclei show comparable translational wave scales only at low temperature.
\paragraph{Ideal gases.}
For number density $n=N/V=P/(k_BT)$, estimate the spacing
\begin{equation}
d\simeq n^{-1/3}=\left(\frac{k_BT}{P}\right)^{1/3}.
\end{equation}
Substituting into $\lambda\gtrsim d$ and raising both sides to appropriate powers gives
\begin{align}
\frac{h^2}{3mk_BT}&\gtrsim\left(\frac{k_BT}{P}\right)^{2/3},\\
T^{5/3}&\lesssim\frac{h^2}{3mk_B}\left(\frac{P}{k_B}\right)^{2/3},
\end{align}
hence
\begin{equation}
\boxed{T\lesssim\frac1{k_B}\left(\frac{h^2}{3m}\right)^{3/5}P^{2/5}.}
\end{equation}
For helium atoms of mass $4.002602\,u$ at $P=101325\,\mathrm{Pa}$, this formula gives
\begin{equation}
\boxed{T_{\rm q,He}\simeq2.94\,\mathrm K.}
\end{equation}
At such temperatures ordinary helium may liquefy, so applying the \emph{ideal gas} estimate at atmospheric pressure is a formal scale comparison rather than a claim of a stable ideal-gas phase.
\paragraph{Hydrogen in interstellar conditions.}
For monatomic hydrogen $m\simeq1.00784u$, $T=3\,\mathrm K$,
\begin{equation}
\lambda=\frac{6.62607015\times10^{-34}}
{\sqrt{3(1.00784u)k_B(3\,\mathrm K)}}
\simeq\boxed{1.45\times10^{-9}\,\mathrm m}.
\end{equation}
Compared with $d\simeq1\,\mathrm{cm}=10^{-2}\,\mathrm m$,
\begin{equation}
\boxed{\frac{\lambda}{d}\simeq1.45\times10^{-7}\ll1}.
\end{equation}
Therefore this extremely dilute \emph{translational} gas is nondegenerate by the stated criterion; the internal electronic structure of its atoms is nevertheless quantum mechanical.
\note{The relevant distance scale differs across problems. A correct conclusion must specify whether the question concerns conduction electrons, atomic centers of mass, or internal bound states.}

\section*{Conclusions}
\addcontentsline{toc}{section}{Conclusions}
Normalization, probabilistic expectations, uncertainty, and probability currents are connected by the Hermitian structure of quantum dynamics. The final exercises contrast this closed-system picture with classical statistical mechanics and phenomenological non-Hermitian evolution.
\begin{thebibliography}{9}
\bibitem{griffiths} D.~J.~Griffiths and D.~F.~Schroeter, \emph{Introduction to Quantum Mechanics}, 3rd ed., Cambridge University Press (2018), Chapter~1, pp.~3--24. DOI: \href{https://doi.org/10.1017/9781316995433}{10.1017/9781316995433}.
\end{thebibliography}
\end{document}
